%=====================================================================
\chapter{Support Vector Machines}\label{ch:svm}
%=====================================================================
\locator{3}{Part III --- Supervised Learning Algorithms}

\begin{outcomes}
By the end of this chapter you will be able to:
\begin{tight}
  \item \textbf{Explain} why the maximum-margin hyperplane is preferred among all
        separating hyperplanes, and \textbf{identify} its support vectors.
  \item \textbf{Compute} the margin of a linear SVM and \textbf{classify} a point
        with its decision function.
  \item \textbf{Describe} the soft-margin SVM and the role of $C$, and
        \textbf{compute} the hinge loss.
  \item \textbf{Explain} the kernel trick, and \textbf{verify} a polynomial kernel
        against its explicit feature map.
  \item \textbf{Tune} $C$ and the RBF width $\gamma$ by cross-validation.
\end{tight}
\end{outcomes}

\begin{prereq}
\chref{math} (dot products, norms), \chref{logreg} (linear decision boundaries,
the loss view of classification) and \chref{generalization} (regularisation and
tuning by cross-validation).
\end{prereq}

\begin{keyterms}
separating hyperplane $\bullet$ margin $\bullet$ maximum-margin classifier
$\bullet$ support vectors $\bullet$ hard margin $\bullet$ soft margin $\bullet$
slack variable $\bullet$ $C$ $\bullet$ hinge loss $\bullet$ dual form $\bullet$
feature map $\bullet$ kernel trick $\bullet$ polynomial kernel $\bullet$ RBF
kernel $\bullet$ $\gamma$ $\bullet$ support vector regression
\end{keyterms}

\section{Which Line?}

When two classes can be separated by a straight line, they can usually be
separated by infinitely many. Logistic regression picks the one that maximises
likelihood; the perceptron of \chref{ann} picks whichever it reaches first. A
\term{support vector machine} picks the line that stays as far as possible from
the nearest examples of both classes --- the one with the widest \term{margin}.
The intuition is that a boundary with room to spare on both sides is the least
likely to be crossed by small perturbations of new data, and the theory agrees:
a wider margin gives a tighter bound on the error on unseen data, regardless of
the number of features.

\section{The Maximum-Margin Hyperplane}

\begin{definitionbox}[Hyperplane, Margin and Support Vectors]
A linear classifier predicts $\operatorname{sign}(\mathbf{w}\cdot\mathbf{x} + b)$
with labels $y \in \{-1, +1\}$. The distance of a point $\mathbf{x}$ from the
hyperplane $\mathbf{w}\cdot\mathbf{x} + b = 0$ is
$|\mathbf{w}\cdot\mathbf{x} + b| / \|\mathbf{w}\|$. Scale $\mathbf{w}$ and $b$ so that
the closest points satisfy $y_i(\mathbf{w}\cdot\mathbf{x}_i + b) = 1$; then the
\term{margin} --- the width of the empty band between the classes --- is
$2/\|\mathbf{w}\|$. The hard-margin SVM solves
\[
\min_{\mathbf{w}, b}\ \tfrac{1}{2}\|\mathbf{w}\|^2
\quad\text{subject to}\quad y_i(\mathbf{w}\cdot\mathbf{x}_i + b) \ge 1 \ \text{ for every } i
\]
The training points that lie exactly on the edges of the band are the
\term{support vectors}. They alone determine the solution: every other point could
be moved or deleted, as long as it stayed outside the band, and the hyperplane
would not change.
\end{definitionbox}

\dsfig{%
\begin{tikzpicture}[font=\scriptsize, scale=0.9,
  ps/.style={circle, fill=secondary, inner sep=2pt},
  ng/.style={rectangle, fill=accent, inner sep=2.2pt}]
% band between f = -1 (x1 + x2 = 0) and f = +1 (x1 + x2 = 4)
\fill[primary!6] (-2,2) -- (2,-2) -- (4.5,-0.5) -- (-0.5,4.5) -- cycle;
\draw[-{Stealth[length=2mm]}, gray!55] (-2.2,0) -- (5,0) node[right, font=\tiny] {$x_1$};
\draw[-{Stealth[length=2mm]}, gray!55] (0,-2.4) -- (0,5) node[above, font=\tiny] {$x_2$};
\draw[primary, line width=1.4pt] (-1.4,3.4) -- (3.4,-1.4);
\draw[primary!70, dashed, line width=0.9pt] (-2,2) -- (2,-2);
\draw[primary!70, dashed, line width=0.9pt] (-0.5,4.5) -- (4.5,-0.5);
\node[ps] at (2,2) {}; \node[ps] at (3,2) {}; \node[ps] at (2,4) {};
\node[ng] at (0,0) {}; \node[ng] at (-1,-1) {}; \node[ng] at (1,-2) {};
\draw[greenhl!70!black, line width=1.2pt] (2,2) circle (0.26);
\draw[greenhl!70!black, line width=1.2pt] (0,0) circle (0.26);
\draw[{Stealth[length=2mm]}-{Stealth[length=2mm]}, goldhl!80!black, line width=0.9pt]
     (0.2,0.2) -- (1.8,1.8);
\node[font=\tiny, goldhl!45!black, anchor=west, fill=white, inner sep=1pt] at (1.25,0.55) {margin $= 2/\|\mathbf{w}\| = 2.83$};
\node[font=\tiny, primary, anchor=west] at (3.45,-1.45) {$f = 0$};
\node[font=\tiny, primary!80, anchor=west] at (4.55,-0.55) {$f = +1$};
\node[font=\tiny, primary!80, anchor=west] at (2.05,-2.05) {$f = -1$};
\node[anchor=north west, align=left, text width=5.4cm, font=\scriptsize, text=gray!55!black]
     at (5.6,4.8)
     {\textbf{The worked example.} Positives (blue) at $(2,2)$, $(3,2)$, $(2,4)$;
      negatives (red) at $(0,0)$, $(-1,-1)$, $(1,-2)$.\\[4pt]
      The widest band touches only $(2,2)$ and $(0,0)$ --- the two
      \textbf{\color{greenhl!55!black}support vectors}, circled. The boundary is the
      line $x_1 + x_2 = 2$, halfway between them.\\[4pt]
      Delete any other point and the answer is unchanged.};
\end{tikzpicture}}%
{The maximum-margin hyperplane. The solid line separates the classes; the dashed lines
bound the widest empty band; the points on the dashed lines are the support
vectors.}{max-margin}

\begin{worked}[A Maximum-Margin Line by Hand]
The six points of \dsref{max-margin}. The closest opposite-class pair is $(2,2)$ and
$(0,0)$, so the boundary is perpendicular to the segment between them:
$\mathbf{w} = a\,(1, 1)$ for some $a > 0$.

\smallskip
\textbf{Canonical scaling.} Put the support vectors exactly on the margins:
\[
\mathbf{w}\cdot(2,2) + b = 4a + b = +1, \qquad \mathbf{w}\cdot(0,0) + b = b = -1
\]
So $b = -1$, $a = 0.5$: $\mathbf{w} = (0.5,\ 0.5)$. The boundary is
$0.5x_1 + 0.5x_2 - 1 = 0$, i.e.\ $x_1 + x_2 = 2$.

\textbf{Margin:} $\|\mathbf{w}\| = \sqrt{0.25 + 0.25} = 0.707$, so
$2/\|\mathbf{w}\| = \mathbf{2.83}$ --- exactly the distance from $(0,0)$ to $(2,2)$.

\textbf{Check the constraints:} $(3,2)$: $f = 1.5 \ge 1$; $(2,4)$: $f = 2 \ge 1$;
$(-1,-1)$: $f = -2 \le -1$; $(1,-2)$: $f = -1.5 \le -1$. All satisfied.

\smallskip
\textbf{Classify} $(1.5,\ 1)$: $f = 0.75 + 0.5 - 1 = 0.25 > 0$, so \textbf{positive}, at
distance $0.25/0.707 = 0.35$ from the boundary --- inside the margin band, so a
point the SVM is not confident about. scikit-learn's \texttt{SVC} returns the same
$\mathbf{w}$, $b$ and support vectors (this chapter's lab).
\end{worked}

\section{Soft Margins and the Hinge Loss}

Real data is rarely separable, and even when it is, one mislabelled point can drag
the hard-margin boundary into a bad place. The \term{soft-margin} SVM lets points
violate the margin at a price. Each point gets a \term{slack variable}
$\xi_i \ge 0$ measuring how far it is on the wrong side of its margin:
\[
\min_{\mathbf{w}, b, \boldsymbol{\xi}}\ \tfrac{1}{2}\|\mathbf{w}\|^2 + C\sum_i \xi_i
\quad\text{subject to}\quad y_i(\mathbf{w}\cdot\mathbf{x}_i + b) \ge 1 - \xi_i
\]
The hyperparameter $C$ sets the price. A large $C$ punishes every violation
heavily: narrow margin, few violations, high variance. A small $C$ tolerates
violations to buy a wider margin: more bias, less variance. $C$ is the SVM's
regularisation knob, and it is tuned by cross-validation.

\begin{definitionbox}[The Hinge Loss]
At the optimum each slack equals $\max(0,\ 1 - y_i f(\mathbf{x}_i))$, so the soft-margin
SVM is equivalent to minimising
\[
\underbrace{\sum_i \max\big(0,\ 1 - y_i f(\mathbf{x}_i)\big)}_{\text{hinge loss}}
\;+\; \frac{1}{2C}\|\mathbf{w}\|^2
\]
--- a loss plus an L2 penalty, the same shape as regularised logistic regression
with the cross-entropy replaced by the \term{hinge}. Points beyond the margin on the
correct side cost nothing.
\end{definitionbox}

\dsfig{%
\begin{tikzpicture}
\begin{axis}[width=8.8cm, height=5.0cm, axis lines=middle,
  xlabel={$y\,f(\mathbf{x})$ (margin of the example)}, ylabel={loss},
  x label style={at={(axis description cs:0.5,-0.08)}, anchor=north},
  y label style={at={(axis description cs:-0.02,0.5)}, rotate=90, anchor=south},
  tick label style={font=\tiny}, label style={font=\scriptsize},
  xmin=-2.2, xmax=3.2, ymin=-0.1, ymax=3.3,
  legend style={font=\tiny, draw=gray!40, at={(0.98,0.97)}, anchor=north east}]
\addplot[gray!70, line width=1.1pt, const plot] coordinates {(-2.2,1)(0,1)(0,0)(3.2,0)};
\addlegendentry{0--1 loss (the error count)}
\addplot[accent, line width=1.4pt, domain=-2.2:3.2, samples=100] {max(0,1-x)};
\addlegendentry{hinge (SVM)}
\addplot[secondary, line width=1.2pt, dashed, domain=-2.2:3.2, samples=100] {ln(1+exp(-x))/ln(2)};
\addlegendentry{logistic (scaled)}
\node[font=\tiny, accent, anchor=south west, align=left] at (axis cs:1.2,0.85)
     {zero beyond the margin:\\these points do not matter};
\end{axis}
\end{tikzpicture}}%
{Three losses as a function of $y f(\mathbf{x})$, positive when the example is
correctly classified. The hinge is exactly zero for points beyond the margin, which is
why only support vectors affect the SVM; the logistic loss is never quite zero.}{hinge}

\begin{worked}[Hinge Losses]
Four training points have margins $y f(\mathbf{x}) = 2,\ 1,\ 0.5,\ -0.5$:
\begin{tight}
  \item $y f = 2$: $\max(0, 1 - 2) = \mathbf{0}$ --- well clear; irrelevant to the solution.
  \item $y f = 1$: $\max(0, 0) = \mathbf{0}$ --- exactly on the margin; a support vector.
  \item $y f = 0.5$: $\mathbf{0.5}$ --- correctly classified but inside the margin.
  \item $y f = -0.5$: $\mathbf{1.5}$ --- misclassified.
\end{tight}
The last three all influence the boundary; the first does not. With $C$ large the
optimiser will work hard to remove the last two losses; with $C$ small it will
accept them for a wider margin.
\end{worked}

\section{The Kernel Trick}

Some problems no straight line can solve --- one class surrounding another, say.
The standard remedy (\chref{linreg}) is to add non-linear features, such as
squares and products, and fit a linear model in the enlarged space. That works,
but the enlarged space can be enormous: all products of pairs of 1{,}000 features
is half a million features.

\dsfig{%
\begin{tikzpicture}[font=\scriptsize,
  ps/.style={circle, fill=secondary, inner sep=1.9pt},
  ng/.style={rectangle, fill=accent, inner sep=2.1pt}]
% 1-D
\draw[-{Stealth[length=2mm]}, gray!60] (-3.3,0) -- (3.4,0) node[right, font=\tiny] {$x$};
\foreach \x in {-2.6,-2.2,-1.8,1.8,2.2,2.7}{\node[ng] at (\x,0) {};}
\foreach \x in {-1.0,-0.5,0.0,0.4,0.9}{\node[ps] at (\x,0) {};}
\node[font=\scriptsize\bfseries, text=primary] at (0,-0.8) {In one dimension: no single threshold separates them};
% 2-D lifted
\begin{scope}[shift={(8.6,-1.0)}]
\draw[-{Stealth[length=2mm]}, gray!60] (-3.3,0) -- (3.4,0) node[right, font=\tiny] {$x$};
\draw[-{Stealth[length=2mm]}, gray!60] (0,-0.2) -- (0,3.6) node[above, font=\tiny] {$x^2$};
\draw[gray!40, domain=-2.9:2.9, samples=40] plot (\x,{0.4*\x*\x});
\foreach \x in {-2.6,-2.2,-1.8,1.8,2.2,2.7}{\node[ng] at (\x,{0.4*\x*\x}) {};}
\foreach \x in {-1.0,-0.5,0.0,0.4,0.9}{\node[ps] at (\x,{0.4*\x*\x}) {};}
\draw[greenhl!70!black, line width=1.3pt] (-3.1,0.85) -- (3.1,0.85);
\node[font=\tiny, greenhl!45!black, anchor=south west] at (1.0,0.9) {a straight line now works};
\node[font=\scriptsize\bfseries, text=primary] at (0,-0.6) {Add the feature $x^2$: separable};
\end{scope}
\draw[-{Stealth[length=3mm]}, gray!60, line width=1.2pt] (3.9,0.3) -- (5.0,0.3)
     node[midway, above, font=\tiny] {$\phi(x) = (x, x^2)$};
\end{tikzpicture}}%
{Lifting data into more dimensions. The inner class (blue) and the outer class (red)
cannot be separated on the line, but after adding $x^2$ as a second feature a horizontal
line separates them; in the original space that line is two thresholds.}{kernel-lift}

\begin{definitionbox}[Kernel]
A \term{kernel} is a function $K(\mathbf{x}, \mathbf{z})$ that equals the dot product
$\phi(\mathbf{x})\cdot\phi(\mathbf{z})$ in some \term{feature space} $\phi$, but is
computed without ever constructing $\phi$. Because the SVM's optimisation and its
decision function use the data only through dot products,
\[
f(\mathbf{x}) = \sum_{i \in \text{SV}} \alpha_i\, y_i\, K(\mathbf{x}_i, \mathbf{x}) + b
\]
replacing every dot product by a kernel fits a linear boundary in the feature space
--- a non-linear boundary in the original one --- at the cost of computing $K$.
This is the \term{kernel trick}. The sum runs over support vectors only: all other
$\alpha_i$ are zero.
\end{definitionbox}

\begin{worked}[The Kernel Computes the Big Dot Product Cheaply]
Take $K(\mathbf{x}, \mathbf{z}) = (\mathbf{x}\cdot\mathbf{z})^2$ in two dimensions, and
$\mathbf{x} = (1, 2)$, $\mathbf{z} = (3, 1)$.

\smallskip
\textbf{The kernel:} $\mathbf{x}\cdot\mathbf{z} = 3 + 2 = 5$, so $K = 25$. One dot product
and one squaring.

\textbf{The long way:} this kernel corresponds to
$\phi(\mathbf{x}) = (x_1^2,\ \sqrt{2}\,x_1 x_2,\ x_2^2)$:
\[
\phi(1,2) = (1,\ 2\sqrt{2},\ 4), \qquad \phi(3,1) = (9,\ 3\sqrt{2},\ 1)
\]
\[
\phi(\mathbf{x})\cdot\phi(\mathbf{z}) = 9 + 2\sqrt{2}\cdot 3\sqrt{2} + 4 = 9 + 12 + 4 = \mathbf{25}
\]
The same number. For degree $p$ on $d$ features the feature space has on the order of
$d^p$ dimensions; the kernel is still one dot product and one power.
\end{worked}

\rowcolors{2}{white}{lightbg}
\begin{longtable}{@{}L{3.0cm}L{4.4cm}L{7.6cm}@{}}
\toprule
\rowcolor{primary!15}
\textbf{Kernel} & \textbf{$K(\mathbf{x}, \mathbf{z})$} & \textbf{Use and settings} \\
\midrule
\endhead
Linear & $\mathbf{x}\cdot\mathbf{z}$ & Many features, text; fast. Tune $C$ only \\
Polynomial & $(\gamma\,\mathbf{x}\cdot\mathbf{z} + r)^p$ & Interactions of up to $p$
features; tune $p$, $C$ \\
RBF (Gaussian) & $\exp(-\gamma\|\mathbf{x} - \mathbf{z}\|^2)$ & The default for
non-linear problems. Large $\gamma$: each point's influence is narrow, the boundary
wiggles (variance). Small $\gamma$: smooth boundary (bias). Tune $C$ and $\gamma$
together \\
\bottomrule
\end{longtable}

\begin{alertbox}[Scale the Features, Then Tune $C$ and $\gamma$ Together]
The RBF kernel is a function of distance, so everything said about scaling in
\chref{knn} applies. And $C$ and $\gamma$ interact --- a larger $\gamma$ often wants a
smaller $C$ --- so search them as a grid, on a logarithmic scale
(\texttt{C} in $\{0.1, 1, 10, 100\}$, \texttt{gamma} in $\{0.001, 0.01, 0.1, 1\}$),
by cross-validation. On the breast cancer data this finds $C = 10$, $\gamma = 0.01$
and 97.9\% accuracy using 64 of the 569 examples as support vectors.
\end{alertbox}

\section{Practicalities}

\textbf{More than two classes.} The SVM is inherently binary. scikit-learn's
\texttt{SVC} trains one classifier per pair of classes (one-vs-one) and lets them
vote; \texttt{LinearSVC} uses one-vs-rest.

\textbf{Probabilities.} The decision function $f(\mathbf{x})$ is a distance, not a
probability. \texttt{SVC(probability=True)} fits a sigmoid to it afterwards
(Platt scaling), at extra cost.

\textbf{Scale of data.} Training cost grows between the square and the cube of the
number of examples, so kernel SVMs suit datasets of up to tens of thousands of
rows. Linear SVMs scale to millions.

\textbf{Regression.} \term{Support vector regression} ignores errors smaller than a
tube of width $\varepsilon$ around the prediction and charges linearly beyond it ---
the regression counterpart of the hinge.

\textbf{Inductive bias.} Among the hypotheses consistent with the data (up to the
slack allowed), prefer the one with the largest margin.

\begin{pitfall}
\begin{tight}
  \item \textbf{Not scaling features.} The RBF kernel and the margin are both
        distances.
  \item \textbf{Tuning $C$ and $\gamma$ separately.} They interact; use a grid.
  \item \textbf{Using a huge $C$ on noisy data.} The hard margin chases every
        mislabelled point.
  \item \textbf{Reading $f(\mathbf{x})$ as a probability.} It is a signed distance;
        calibrate if probabilities are needed.
  \item \textbf{Running a kernel SVM on a million rows.} Use a linear SVM, or a
        different model.
\end{tight}
\end{pitfall}

%---------------------------------------------------------------------
\begin{fourlenses}{Support Vector Machines}
\lensLA{With a few hundred students and a dozen engagement features, an RBF SVM is
often among the most accurate classifiers available --- and among the hardest to
explain. The support vectors do give one explanation: the borderline students the
model learned the boundary from, which advisers can examine directly.}

\lensPM{Text classification of requirement statements (functional or not,
ambiguous or not) with a linear SVM over word counts is a classic, strong
baseline: high-dimensional, sparse data is where linear SVMs are at their best.}

\lensIOT{A trained linear SVM is a single weight vector, so inference on a device
is one dot product. Kernel SVMs must store their support vectors --- often
thousands --- which may not fit; count them before choosing.}

\lensSEC{SVMs were for years the standard malware and intrusion classifier over
static file features. The margin gives some robustness to small changes an attacker
makes, but only in the directions the features capture; an attacker who changes
features the model ignores is not slowed at all.}
\end{fourlenses}

%---------------------------------------------------------------------
\begin{lab}[Margins, Kernels and a Tuning Grid]
\begin{lstlisting}[language=Python]
import numpy as np
from sklearn.svm import SVC
from sklearn.datasets import make_circles, load_breast_cancer
from sklearn.model_selection import cross_val_score, GridSearchCV
from sklearn.pipeline import make_pipeline
from sklearn.preprocessing import StandardScaler

# --- 1. the worked example: a maximum-margin line ------------------------
X = np.array([[2, 2], [3, 2], [2, 4], [0, 0], [-1, -1], [1, -2]])
y = np.array([1, 1, 1, -1, -1, -1])
svm = SVC(kernel="linear", C=1e6).fit(X, y)       # huge C = hard margin
w, b = svm.coef_[0], svm.intercept_[0]
print("w =", w.round(3), " b =", round(b, 3))
print("margin width 2/||w|| =", round(2 / np.linalg.norm(w), 3))
print("support vectors:", svm.support_vectors_.tolist())
print("f(1.5, 1) =", round(float(svm.decision_function([[1.5, 1]])[0]), 3))

# --- 2. a problem no line can solve --------------------------------------
Xc, yc = make_circles(n_samples=400, noise=0.08, factor=0.4, random_state=0)
for kernel in ("linear", "poly", "rbf"):
    acc = cross_val_score(SVC(kernel=kernel, degree=2), Xc, yc, cv=5).mean()
    print(f"circles, {kernel:6s} kernel: accuracy {acc:.3f}")

# --- 3. tune C and gamma together, features scaled -----------------------
Xb, yb = load_breast_cancer(return_X_y=True)
pipe = make_pipeline(StandardScaler(), SVC(kernel="rbf"))
grid = GridSearchCV(pipe, {"svc__C": [0.1, 1, 10, 100],
                           "svc__gamma": [0.001, 0.01, 0.1, 1]}, cv=5)
grid.fit(Xb, yb)
print("best:", grid.best_params_, "CV accuracy",
      round(grid.best_score_, 3))
print("support vectors used:", grid.best_estimator_[-1].n_support_.sum(),
      "of", len(yb))
\end{lstlisting}
In part 1, add the point $(1, 1.5)$ with label $+1$ and rerun. It lies inside the old
margin: which support vectors change, and what happens to the margin width?
\end{lab}

%---------------------------------------------------------------------
\begin{checkpoint}
\begin{tightnum}
  \item A linear SVM has $\mathbf{w} = (3, 4)$ and $b = -5$. What is its margin width,
        and how is $(1, 1)$ classified?
  \item What happens to the margin and to training errors as $C$ is increased?
  \item Why can the kernel trick be applied to the SVM at all?
\end{tightnum}
\end{checkpoint}

%---------------------------------------------------------------------
\begin{chaptersummary}
\begin{tight}
  \item Of all separating hyperplanes the SVM chooses the one with the widest
        margin, $2/\|\mathbf{w}\|$; it is determined by the support vectors alone.
  \item The soft-margin SVM allows violations priced by $C$, and is equivalent to
        hinge loss plus an L2 penalty. Large $C$: narrow margin, high variance.
  \item The hinge loss is zero beyond the margin, which is why only support vectors
        matter.
  \item Kernels compute dot products in a large feature space without building it;
        the RBF kernel is the default for non-linear problems.
  \item Scale the features and tune $C$ and $\gamma$ together on a logarithmic grid
        by cross-validation.
\end{tight}
\end{chaptersummary}

\begin{reviewq}
  \item \textbf{(MCQ)} The support vectors of an SVM are: (a)~all training points
        (b)~the points on or inside the margin (c)~the misclassified points only
        (d)~the cluster centres.
  \item \textbf{(MCQ)} With the RBF kernel, a very large $\gamma$ tends to:
        (a)~underfit (b)~overfit (c)~have no effect (d)~make the model linear.
  \item \textbf{(MCQ)} The hinge loss of a point with $y f(\mathbf{x}) = 0.3$ is:
        (a)~0 (b)~0.3 (c)~0.7 (d)~1.3.
  \item \textbf{(Short)} Why is a wider margin expected to generalise better?
  \item \textbf{(Short)} Compare the hinge loss with the logistic loss.
  \item \textbf{(Short)} Explain the kernel trick without formulas.
  \item \textbf{(Applied)} For support vectors $(1, 3)$ (label $+1$) and $(1, 1)$
        (label $-1$), find $\mathbf{w}$, $b$ and the margin of the maximum-margin
        line, and classify $(4, 2.5)$.
  \item \textbf{(Applied)} Verify that $K(\mathbf{x}, \mathbf{z}) = (\mathbf{x}\cdot\mathbf{z})^2$
        equals $\phi(\mathbf{x})\cdot\phi(\mathbf{z})$ for $\mathbf{x} = (2, -1)$ and
        $\mathbf{z} = (1, 3)$.
\end{reviewq}
